\section{Kiểm thử và đánh giá}

Phần kiểm thử tập trung đánh giá độ chính xác của thuật toán rule-based trên
tập dữ liệu gốc \path{data/raw/pii_dataset.csv}. Script kiểm thử được đặt tại
\path{tests/test_rule_based_accuracy.py}. Với mỗi dòng dữ liệu, hệ thống
lấy trường \texttt{text} làm đầu vào, chạy pipeline phát hiện thực thể hiện tại
gồm \texttt{regex\_detector}, \texttt{context\_detector} và \texttt{merger}, sau
đó so sánh kết quả nhận diện với các cột ground truth trong file CSV.

Các loại thực thể được đánh giá trong phần này gồm:

\begin{center}
    \texttt{NAME}, \texttt{EMAIL}, \texttt{PHONE}, \texttt{URL}
\end{center}

Trong đó, \texttt{NAME} được so sánh với cột \texttt{name}, \texttt{EMAIL} với
cột \texttt{email}, \texttt{PHONE} với cột \texttt{phone}, và \texttt{URL} với
cột \texttt{url}. Đối với số điện thoại, script chuẩn hóa về chuỗi chữ số trước
khi so sánh nhằm tránh sai lệch do dấu cách, dấu ngoặc hoặc dấu gạch ngang.

\subsection{Các chỉ số đánh giá}

Script in ra các chỉ số sau:

\begin{itemize}[nosep]
    \item \textbf{Support}: số dòng có giá trị ground truth cho loại thực thể
          tương ứng.
    \item \textbf{Correct}: số dòng mà detector tìm được đúng giá trị ground
          truth trong danh sách dự đoán.
    \item \textbf{Accuracy}: tỉ lệ \texttt{Correct / Support}. Chỉ số này cho
          biết trong các dòng có ground truth, hệ thống có tìm được đúng thực
          thể cần tìm hay không.
    \item \textbf{Precision}: tỉ lệ \texttt{TP / (TP + FP)}. Chỉ số này cho
          biết trong các thực thể hệ thống dự đoán ra, bao nhiêu thực thể là
          đúng.
    \item \textbf{Recall}: tỉ lệ \texttt{TP / (TP + FN)}, tương đương khả năng
          tìm lại đúng các thực thể có trong ground truth.
    \item \textbf{F1}: trung bình điều hòa giữa precision và recall.
\end{itemize}

Vì mỗi dòng CSV chỉ có một giá trị ground truth cho từng cột như \texttt{name}
hoặc \texttt{email}, \texttt{Accuracy} trong bảng dưới đây được hiểu là tỉ lệ
tìm đúng giá trị chính của dòng đó. Chỉ số này khác với \texttt{Precision}: một
detector có thể tìm đúng tên thật trong dòng văn bản nhưng đồng thời dự đoán dư
nhiều cụm khác là \texttt{NAME}; khi đó accuracy vẫn cao nhưng precision thấp.

\subsection{Kết quả thực nghiệm}

Kết quả được ghi nhận khi chạy lệnh:

\begin{verbatim}
python tests/test_rule_based_accuracy.py
\end{verbatim}

Tổng số dòng được đánh giá là 4434. Kết quả chi tiết được trình bày trong
Bảng~\ref{tab:rule_based_accuracy}.

\begin{table}[H]
    \centering
    \caption{Kết quả đánh giá rule-based detector trên tập \texttt{pii\_dataset.csv}}
    \label{tab:rule_based_accuracy}
    \renewcommand{\arraystretch}{1.25}
    \begin{adjustbox}{max width=\textwidth}
        \begin{tabular}{lrrrrrrrr}
            \toprule
            \textbf{Entity}  & \textbf{Support}   & \textbf{Correct} & \textbf{Accuracy}
                             & \textbf{Precision} & \textbf{Recall}  & \textbf{F1}
                             & \textbf{FP}        & \textbf{FN}                                                                       \\
            \midrule
            \texttt{NAME}    & 4434               & 3845             & 86.72\%           & 18.39\% & 86.72\% & 30.35\% & 17060 & 589  \\
            \texttt{EMAIL}   & 4434               & 3825             & 86.27\%           & 95.79\% & 86.27\% & 90.78\% & 168   & 609  \\
            \texttt{PHONE}   & 4434               & 2383             & 53.74\%           & 99.04\% & 53.74\% & 69.68\% & 23    & 2051 \\
            \texttt{URL}     & 1322               & 526              & 39.79\%           & 86.23\% & 39.79\% & 54.45\% & 84    & 796  \\
            \midrule
            \textbf{Overall} & 14624              & 10579            & 72.34\%           & 37.90\% & 72.34\% & 49.74\% & 17335 & 4045 \\
            \bottomrule
        \end{tabular}
    \end{adjustbox}
\end{table}

\subsection{Nhận xét}

Kết quả cho thấy nhóm thực thể có cấu trúc rõ ràng như \texttt{EMAIL} đạt hiệu
quả tốt nhất. Precision của \texttt{EMAIL} đạt 95.79\%, cho thấy phần lớn email
được detector dự đoán ra là chính xác. Tuy nhiên recall chỉ đạt 86.27\%, nghĩa
là vẫn còn một số email trong dataset chưa được phát hiện, thường do khác biệt
format hoặc dấu câu đi kèm trong văn bản.

Đối với \texttt{PHONE}, precision rất cao, đạt 99.04\%, nhưng recall chỉ đạt
53.74\%. Điều này cho thấy regex số điện thoại khá thận trọng: khi đã nhận diện
thì thường đúng, nhưng bỏ sót nhiều trường hợp có định dạng đa dạng. Đây là một
điểm cần cải thiện nếu muốn tăng độ phủ của hệ thống.

\texttt{URL} có precision tương đối cao, đạt 86.23\%, nhưng accuracy và recall
chỉ đạt 39.79\%. Nguyên nhân chính là dataset có nhiều URL hoặc placeholder URL
không theo đúng dạng mà regex hiện tại ưu tiên, chẳng hạn các dạng thiếu scheme
hoặc bị gắn với dấu câu trong câu.

Với \texttt{NAME}, accuracy đạt 86.72\%, nghĩa là hệ thống thường tìm được tên
chính của dòng dữ liệu. Tuy nhiên precision chỉ đạt 18.39\% do detector theo
ngữ cảnh còn dự đoán dư nhiều cụm viết hoa không phải tên người. Điều này phản
ánh đặc thù của rule-based name detection: nếu rule rộng thì recall cao hơn
nhưng dễ phát sinh false positive; nếu rule chặt hơn thì precision tăng nhưng
có nguy cơ bỏ sót tên thật.

Nhìn chung, kết quả overall đạt accuracy 72.34\% và F1 49.74\%. Hệ thống
rule-based hiện tại phù hợp với các thực thể có cấu trúc cố định, đặc biệt là
\texttt{EMAIL} và các số điện thoại có format rõ ràng. Các thực thể phụ thuộc
nhiều vào ngữ cảnh hoặc có format đa dạng như \texttt{NAME} và \texttt{URL} cần
được bổ sung thêm rule, dictionary hoặc cơ chế lọc false positive để cải thiện
precision và recall.
